\documentclass[a4paper,12pt]{memoir}
\usepackage{graphicx}
\usepackage[usenames,svgnames]{xcolor}
\chapterstyle{bianchi}
\usepackage{a4wide,amsfonts,amsthm,amsmath,appendix}
\usepackage[show]{ed}
\usepackage{url}
\usepackage[colorlinks=true,linkcolor=blue,urlcolor=blue]{hyperref}
\usepackage{xspace}

% Less detailed TOC
\setcounter{tocdepth}{1}
\def\arxiv{\textsc{arXiv}\xspace}

\newcommand{\labelentry}{.}
\newcounter{entryi}


%%% ENTRIES for expression case study:
\newenvironment{entries}[2]%
{\begin{table}\def\capentries{#1}\def\labelentries{#2}%
\setcounter{entryi}{1}\begin{tabular}{|llll|}\hline  & Expression & Meaning & Syntax \\ \hline}%
{\hline\end{tabular}\caption{\capentries}\label{\labelentries}\end{table}}%

\newcommand\entry[4]{{\theentryi\labelentry}\stepcounter{entryi} & #1 & #2 & #3 \\[1mm] & \textbf{Discussion:}& \multicolumn{2}{l|}{#4} \\[3mm]}

%DECLARATION SECTION
\theoremstyle{plain}
\newtheorem{thm}[equation]{Theorem}
\newtheorem{lem}[equation]{Lemma}
\newtheorem{cor}[equation]{Corollary}
\newtheorem{prop}[equation]{Proposition}

\theoremstyle{definition}
\newtheorem{defi}[equation]{Definition}
\newtheorem{rmk}[equation]{Remark}
\newtheorem{example}[equation]{Example}

\theoremstyle{remark}
\newtheorem{notation}{Notation}

%COMMAND SECTION
%\overfullrule=5pt
\renewcommand{\textwidth}{390pt}
\renewcommand{\baselinestretch}{1.5}
\numberwithin{equation}{section}

\newcommand{\bib}{\bibitem}
\newcommand{\ra}{\rightarrow}
\newcommand{\lra}{\longrightarrow}
\newcommand{\nin}{\noindent}
\newcommand{\bsk}{\bigskip}
\newcommand{\msk}{\medskip}
\newcommand{\ssk}{\smallskip}

\newcommand{\contrac}{\mathop{\raise.45ex\hbox{$\underline{\hskip7pt}$}\raise.5ex\hbox{$\mkern-2mu\scriptstyle|$}}\nolimits}

%Greek letters
\newcommand{\al}{\alpha}
\newcommand{\be}{\beta}
\newcommand{\ga}{\gamma}
\newcommand{\Ga}{\Gamma}
\newcommand{\de}{\delta}
\newcommand{\De}{\Delta}
\newcommand{\la}{\lambda}
\newcommand{\La}{\Lambda}
\newcommand{\Th}{\Theta}
\newcommand{\si}{\sigma}
\newcommand{\Si}{\Sigma}
\newcommand{\om}{\omega}
\newcommand{\Om}{\Omega}
\newcommand{\ka}{\kappa}
\newcommand{\ep}{\epsilon}
\newcommand{\vep}{\varepsilon}
\newcommand{\vphi}{\varphi}
\newcommand{\na}{\nabla}


%Caligraphic
\newcommand{\CA}{\mathcal A}
\newcommand{\CB}{\mathcal B}
\newcommand{\CC}{\mathcal C}
\newcommand{\CCD}{\mathcal D}
\newcommand{\CE}{\mathcal E}
\newcommand{\CF}{\mathcal F}
\newcommand{\CG}{\mathcal G}
\newcommand{\CH}{\mathcal H}
\newcommand{\CI}{\mathcal I}
\newcommand{\CJ}{\mathcal J}
\newcommand{\CK}{\mathcal K}
\newcommand{\CL}{\mathcal L}
\newcommand{\CM}{\mathcal M}
\newcommand{\CN}{\mathcal N}
\newcommand{\CO}{\mathcal O}
\newcommand{\CP}{\mathcal P}
\newcommand{\CQ}{\mathcal Q}
\newcommand{\CR}{\mathcal R}
\newcommand{\CS}{\mathcal S}
\newcommand{\CT}{\mathcal T}
\newcommand{\CU}{\mathcal U}
\newcommand{\CV}{\mathcal V}
\newcommand{\CW}{\mathcal W}
\newcommand{\CX}{\mathcal X}
\newcommand{\CY}{\mathcal Y}
\newcommand{\CZ}{\mathcal Z}

%Blackboard bold
\newcommand{\BP}{\mathbb P}
\newcommand{\BT}{\mathbb T}
\newcommand{\BR}{\mathbb R}
\newcommand{\BZ}{\mathbb Z}
\newcommand{\BC}{\mathbb C}
\newcommand{\BV}{\mathbb V}
\newcommand{\BQ}{\mathbb Q}

%Text
\newcommand{\te}{\text{e}}
\newcommand{\ti}{\text{i}}

%Bold letters
\newcommand{\bq}{\mathbf{q}}
\newcommand{\bl}{\mathbf{l}}
\newcommand{\bx}{\mathbf{x}}
\newcommand{\bw}{\mathbf{w}}
\newcommand{\bm}{\mathbf{m}}
\newcommand{\bM}{\mathbf{M}}
\newcommand{\bA}{\mathbf{A}}

%Commands specific for file "torus" 
\newcommand{\CLh}{\hat{\mathcal L}}
\newcommand{\CLP}{\mathcal{L}_{\mathcal{P}}}
\newcommand{\CHP}{\mathcal{H}_{\mathcal{P}}}
\newcommand{\deph}{ {\left| \text{Det} \,{\mathcal P} \right|}^{\frac{1}{2}} }
\newcommand{\dephs}{ {\left| \text{Det} \,{\mathcal P}^{*} \right|}^{\frac{1}{2}} }
\newcommand{\lag}{Lag (\mathcal V)}
\newcommand{\tlag}{\widetilde{Lag}(\mathcal V)}
\newcommand{\flag}{Lag_4 (\mathcal V)}


\begin{document}
%% Some shades
\definecolor{Dark}{gray}{0.2}
\definecolor{MedDark}{gray}{0.4}
\definecolor{Medium}{gray}{0.6}
\definecolor{Light}{gray}{0.8}
\newlength{\drop}% for my convenience
\newlength{\dropt}% for my convenience

\renewcommand{\maketitle}{\begingroup%
\drop = 0.1\textheight
\dropt = 0.05\textheight
\fboxsep 0.5\baselineskip
\sffamily
\vspace*{\dropt}
\centering
{\textcolor{NavyBlue}{\HUGE The Structure of Mathematical\\[3mm] Expressions}}\par
\vspace{0.5\drop}
\fcolorbox{NavyBlue}{Dark}{\textcolor{white}{\normalfont\Large
An {\textsc{arXiv}} Case Study}}\par
\vspace{\drop}
{\Large Deyan Ginev and Bruce R. Miller}\par
\vspace{0.05\textheight}
{\large National Institute of Standards and Techonology}\par
\vfill
{\begin{center}\today\end{center}}
\vspace*{\drop}
{\hfill\includegraphics{NIST}}
\endgroup}
\makeatother

\pagestyle{empty}
\maketitle
\newpage
\pagestyle{plain}
% Remove parskip for toc
\setlength{\parskip}{0ex plus 0.5ex minus 0.2ex}
\tableofcontents

\chapter{Introduction}
In this study, we survey the notational diversity of present-day mathematical expressions, in order to uncover their linguistic phenomena. A practical motivation for this study is to provide a foundation for determining the boundary between syntactic and semantic phenomena in said expressions, from the perspective of language modeling. The ultimate goal of this project is to construct a grammar of mathematical expressions, which captures all relevant syntactic properties established in this study, and allows for the semantic analysis necessary to model and observe the semantic relationships.

\section{Motivation}
We want to enable machine-reading of formulas, in order to provide a variety of user-assistance services, such as semantic search, text-to-speech synthesis, semantic interactions (definition lookup), as well as computer algebra support (``evaluate subexpressions on demand'') and ultimately computer verification (``does that proof step really hold?'').\ednote{expand}
\section{Related Resources}
Notation census, beginnings of study are in Deyan's thesis, Naproche and FMathL have examples, but no real systematic study.\ednote{expand}

\chapter{Methods}
\section{Training Corpus}
The primary corpus on which we base this investigation is the Cornel pre-print archive ``\arxiv''\ednote{cite here}, consisting of over 700,000 articles in 37 scientific subfields.
\subsection{\arxiv Sandbox}
\ednote{Say that, on the \arxiv front, we first start with the train sandbox from Deyan's thesis}
\begin{table}\begin{center}
\begin{tabular}{|ll|}\hline
Train1 & Differential Geometry \\ & \url{http://arxmliv.kwarc.info/files/9609/dg-ga.9609012} \\[2mm]
Train2 & Quantum Physics \\  & \url{http://arxmliv.kwarc.info/files/0910/0910.5733/} \\[2mm]
Train3 & High Energy Physics - Theory \\  & \url{http://arxmliv.kwarc.info/files/9407/hep-th.9407125/} \\[2mm]
Train4 & Commutative Algebra \\  & \url{http://arxmliv.kwarc.info/files/0809/0809.4873/} \\[2mm]
Train5 & Statistics Theory \\  & \url{http://arxmliv.kwarc.info/files/0905/0905.1486/} \\[2mm]
Train6 & General Relativity and Quantum Cosmology \\  & \url{http://arxmliv.kwarc.info/files/0807/0807.2507/} \\[2mm]
Train7 & Cosmology and Extragalactic Astrophysics \\  & \url{http://arxmliv.kwarc.info/files/0908/0908.2548} \\[2mm]
Train8 & Exactly Solvable and Integrable Systems \\  & \url{http://arxmliv.kwarc.info/files/0905/0905.2033} \\[2mm]
Train9 & Geometric Topology \\  & \url{http://arxmliv.kwarc.info/files/0809/0809.4477} \\[2mm]
Train10 & Algebraic Geometry \\  & \url{http://arxmliv.kwarc.info/files/0704/0704.0537} \\ \hline
\end{tabular}
\caption{Sandbox of Ten Random \arxiv Papers from Diverse Scientific Subfields}
\end{center}
\end{table}

As a secondary resource, we we will also consult entry-level literature on highschool mathematics, in order to exhibit basic phenomena, as well as to demonstrate phenomena apriori known to the authors.\ednote{Wikipedia? PEMDAS?}

\section{Structural Annotation}

As one of the goals of our study is to establish a first guess of an underspecified operator tree\ednote{make sure the concepts are introduced and/or rephrase}, any annotation must at its core mark up the applicative logical structure of the mathematical expression. This process will build up a formula tree, the collection of which can later be used as a gold standard for developing a grammatical model of the language of symbolic mathematics.

\section{Annotation Vocabulary}

Another core goal is to discover and describe interesting linguistic phenomena that occur naturally in our corpus. Examples of what we consider ``interesting'' are phenomena that induce ambiguity, or legitimize what would typically be ungrammatical fragments. Cases of ambiguity are well-known to follow from semantic overloading of symbols, implicit argument scopes of operations or eliding syntax, leaving the reader with the task of guessing the ``invisible'' dynamics.Use of custom shorthands, however, as well as custom notations in general, expands the grammar of symbolic mathematics, often in completely non-standard ways that can only be grasped through a deep understanding of the document at hand.

As multiple interesting observations can be made for a single large mathematical formula, it is natural to annotate multiple relevant subexpressions. More concretely, for each phenomenon of interest, we annotate the greatest common subtree (GCT) of all participating subtrees. In case we find a long-range relationship in a large formula, the annotation would hence be placed on the formula root.

The annotations can be utilized for different purposes - browsing by specific phenomena, syntactic feature or lemma, training a classifier, etc. Thus, we take a compositional, standardized approach to providing labels from a fixed vocabulary for the relevant ontological classes of structural properties.

\begin{table}
\begin{center}
\begin{tabular}{|l|l|} 
\hline \textbf{Property} & \textbf{Keywords} \\ \hline
\textbf{Fixity} & prefix, infix, postfix, superfix, subfix, circumfix, transfix, nofix\footnote{``atom'' may also be known as ``nofix''} \\[1mm]
\textbf{Role (Symbols)} & separator, modifier, relation,  operator, metarelation \\[1mm]
\textbf{Role (Objects)} & factor, term, statement, variable, constant, modified \\[1mm]
\textbf{Role (Structure)} & tuple, sequence, expression, shorthand, notation \\[1mm]
\textbf{Composition} & invisible, atom, complex, chained\\[1mm]
\textbf{Shallow Semantics} & type, function, constructor, other \\[1mm]
\textbf{Linguistic} & ellipsis, metonymy, ambiguity, vagueness, anaphora\\[1mm]
\textbf{Math Practices} & framing\\[1mm]
\hline
\end{tabular}
\caption{Keyword Vocabulary for Syntactic Properties}
\end{center}
\end{table}


\chapter{A Study of Mathematical Syntax}

\section{Basics}
\subsection{Foundations}

\ednote{arithmetic, grouping fences and equality}
\ednote{basic relations and orderings}
\ednote{arithmetic and algebraic sequences?} 

% Maybe add a high school subsection/subsubsection:
\subsection{High School}
\ednote{geometry here, otherwise a separate geometry subsection}
\ednote{trigonometry, complex and rational numbers}

\section{Discrete math}
\subsection{Set Theoretic Notations}
\ednote{elementhood, inclusions, set constructors, overloaded arith ops}
\ednote{also maps : domains -> codomains, xRy notations}

\begin{entries}{Set Theory Notations, Part 1}{study:sets1}
\entry{$ W \in \CP \cap \CZ $}{set membership}%
{ \MathTree [.\is{infix,relation}{$\in$} W [.$\cap$ $\CP$ $\CZ$ ] ] }%
{set ops precede set relations, [Train1]}
% \entry{$ \; \nu : \overset{n}{\times} \BV \ra \BR$}{a map}{typed modifier}{$n$-ary cross-product, [Train1]}
% \entry{\footnotesize $\CZ^* = \{ X \in \CV \mid \om(X,W) \in \BZ, \text{ for all }  W \in \CZ \}$}{definition to set}{typed modifier}{NL mixins, quantified relation, [Train1]}
% \entry{$\text{span}_{\BR} \{ W_{1} , \dots ,W_{g} \} $}{span of a set}{prefix operator modified}{set operators can take fenced yet not simply \textit{grouped} arguments, [Train1]}

% \entry{$\mathcal{BS}_{\CP} = \underset{\bq \in (\BZ/ k \BZ)^g}{\cup} \, \La_{\bq}$}{definitional assignment}{infix relation}{$n$-ary union, ranges over subscript, [Train1]}
% \entry{ $L_{12} = L_{\CP_1} \cap L_{\CP_2} \neq 0$}{non-transversal intersection}{modified chained}{acts as relation, is it modified or chained?, [Train1]}
% \entry{ $\al \in GL(g,\BZ) \subset Sp(2g,\BZ)$}{statement}{infix relation}{membership comes after set inclusion, due to well-typedness!, [Train1]}
\end{entries}

%\Tree [.\is{infix,relation}{$\in$} [W .$\cap$ [$\CP$ $\CZ$] ] ]
\subsection{Logical Operators}
\ednote{classic logic, HOL, type theories} 
\subsection{Combinatorics}
\ednote{Infinite sums}
\ednote{binomials, combinations, permutations,}
\subsection{Number Theory}
\ednote{modulo modifiers}
\ednote{tuples}
\ednote{divisibility notations $a \mid b $ and $b/a$}
\ednote{DLMF sneaky notations}
\subsection{Graph Theory}
\ednote{edge and vertex notations}
\ednote{incidence and adjacency notations}
\ednote{Wiki is very nice: \url{http://en.wikipedia.org/wiki/Glossary\_of\_graph\_theory}}
\subsection{Algebra}
\ednote{vectors}
\ednote{maps and complements}
\ednote{groups}
\ednote{lattices}
\documentclass{article}
\usepackage{latexml}
\lxDeclare[role=ID]{$a$}%
\lxDeclare[role=ID]{$b$}%
\lxDeclare[role=ID]{$c$}%
\lxDeclare[role=ID]{$d$}%
\lxDeclare[role=ID]{$e$}%
\lxDeclare[role=FUNCTION]{$f$}%
\lxDeclare[role=FUNCTION]{$g$}%
\lxDeclare[role=FUNCTION]{$h$}%
\lxDeclare[role=OPFUNCTION]{$F$}%
\lxDeclare[role=OPFUNCTION]{$G$}%
\lxDeclare[role=OPFUNCTION]{$H$}%
\begin{document}
\section{Basic Functions}
\[ g(a) = g(b) \]
\[ f(a) = g(h(c)) \]
\[ h(c) = \sin a^2 \]
\[ \sin f(a) \]

\section{Trigonometric functions}
\[ \sin a + \tan b \]
\[ \sin a \cos b \neq \sin(a\cos b)\]
% Parses wrong ? but how to clarify?
\[ \sin 2\pi \neq \log 2x^2 \]

\section{Missing arguments}
% Presumably these are products
\[ f g h \]
\[ f g h a \]
\section{Operator-like functions}
% But these are applied!
\[ F G H \]
\[ F G H a \]

\[ f'(a) = a^2 \]
\[ f^2(a) = g^b(2) \]

\section{Scripts}
\[ \ln a = \log_e a = \ln_{+} b\]
\end{document}


\section{Continuous math}
\subsection{Calculus}
\ednote{differentials, integrals, limits, remember brownian motion integral notations!}
\subsection{Probability}

\ednote{Bayes formula with multiple denotations of P} % syntactically irrelevant?
\ednote{Various conditional and joint probability notations}

\subsection{Interval Notation and Arithmetic}
\ednote{introduce interval notations, then move to interval arithmetic}
\subsection{Topology}
\ednote{manifold constructors and notations}
%Commands specific for file "torus" 
\newcommand{\CLh}{\hat{\mathcal L}}
\newcommand{\CLP}{\mathcal{L}_{\mathcal{P}}}
\newcommand{\CHP}{\mathcal{H}_{\mathcal{P}}}
\newcommand{\deph}{ {\left| \text{Det} \,{\mathcal P} \right|}^{\frac{1}{2}} }
\newcommand{\dephs}{ {\left| \text{Det} \,{\mathcal P}^{*} \right|}^{\frac{1}{2}} }
\newcommand{\lag}{Lag (\mathcal V)}
\newcommand{\tlag}{\widetilde{Lag}(\mathcal V)}
\newcommand{\flag}{Lag_4 (\mathcal V)}

\subsection{Differential Geometry}
Some intro text?\\

\begin{entries}{Differential Geometry Notations, Part 1}{study:diffgeometry1}
\entry{$(\CV/ \CZ, k \om)$}{symplectic torus}{circumfix constructor}{[Train1]}
\entry{$\CZ$}{self-dual lattice}{atom abbreviation}{[Train1]}
\entry{$(\CV, \om)$}{symplectic vector space}{circumfix constructor}{[Train1]}
\entry{$\lag$}{Lagrangian Grassmannian}{circumfix constructor}{[Train1]}
\entry{$\flag$}{4-fold covering space}{applicative constructor}{[Train1]}
\entry{$\CM_{\Si}$}{moduli space}{scripted constructor}{[Train1]}
\entry{$\Si$}{Riemann surface}{atom variable}{[Train1]}
\entry{$H^1(\Si;\BR)$}{chomology space}{applicative constructor}{[Train1]}
\entry{$H^1(\Si;\BR)/ H^1(\Si;\BZ)$}{torus}{applicative constructor}{[Train1]}
\entry{$(M , \om)$}{symplectic manifold}{circumfix constructor}{[Train1]}
\entry{$f \in \CC^{\infty} (M)$}{smooth function}{modified atom}{[Train1]}
\entry{$ X_{f} $}{field}{scripted constructor}{[Train1]}
\entry{$\contrac$}{interior product}{complex infix operator}{Formed via $\backslash$mathop in {\TeX}, [Train1]}
\entry{$[ \om ] \in H^{2}(M; \BR)$}{cohomology class}{modified complex object}{[Train1]}
\entry{$( \cdot , \cdot )$ }{hermitian metric}{notation pattern}{[Train1]}
\end{entries}

\begin{nextentries}{Differential Geometry Notations, Part 2}{study:diffgeometry2}
\entry{$-2 \pi \ti \, \om$}{curvature of $\na$}{arithmetic expression}{[Train1]}
\entry{$(\CL , \na)$}{prequantum line bundle}{circumfix constructor}{[Train1]}
\entry{$U \subset M$}{open subset}{modified atom}{[Train1]}
\entry{$\CL \bigr|_{U}$}{restricted line bundle}{modified atom}{postfix restriction via ``$\bigr|_{U}$'', [Train1]}
\entry{$s \in \Ga(U ;\CL)$}{nonzero section}{modified atom}{[Train1]}
\entry{$\na s = - 2 \pi \ti \, \theta \, s$}{equation}{relation}{[Train1]}
\entry{$\om \bigr|_{U} = d \theta$}{equation}{relation}{[Train1]}
\entry{$T_{x} M$}{bundle}{applicative constructor}{invisible infix bundle-forming operator, [Train1]}
\entry{$\om \bigr|_{\CP_{x}} \equiv 0$}{equivalence}{relation}{[Train1]}
\entry{$\dim \CP_{x} = \frac{1}{2} \, \dim T_{x} M$}{equality}{relation}{$\dim$ has lower precedence than invisible bundle-formation, [Train1]}
\entry{$ [X,Y] \in \CX_{\CP}(M)$}{commutator is in set}{relation}{used as verb phrase in sentence, [Train1]}
\entry{$\na^{\CP}$}{covariant differentiation}{scripted prefix op}{big op?, [Train1]}
\entry{\parbox{3cm}{\scriptsize\begin{alignat*}{3}
  \na^{\CP} :  \CX_{\CP}(M)  &\times \CX_{\CP}(M) & \lra &  \;   \CX_{\CP}(M) \\
     ( X  & , Y ) & \longmapsto & \; \na^{\CP}_X \, Y  \, ,\notag
\end{alignat*}}}{domain specification}{typing modifier}{alignment splits type statement, trailing comma [Train1]}
\entry{$(\na^{\CP}_X \, Y) \, \contrac \, \om = X \, \contrac \, d \, (Y \, \contrac \, \om) \, .$}%
{definitional assignment}{relation}{trailing dot, [Train1]}
\entry{$\Pi_{\CP} : M \ra M/ \CP$}{canonical projection map}{typed modifier}{[Train1]}
\end{nextentries}

% % Extension template
\begin{nextentries}{Differential Geometry Notations, Part 3}{study:diffgeometry3}
\entry{$T^g$}{$g$-dimensional torus}{complex object}{script means dimensionality[Train1]}
\entry{}{}{}{[Train1]}
\entry{}{}{}{[Train1]}
\entry{}{}{}{[Train1]}
\entry{}{}{}{[Train1]}
\entry{}{}{}{[Train1]}
\entry{}{}{}{[Train1]}
\entry{}{}{}{[Train1]}
\entry{}{}{}{[Train1]}
\entry{}{}{}{[Train1]}
\entry{}{}{}{[Train1]}
\entry{}{}{}{[Train1]}
\entry{}{}{}{[Train1]}
\entry{}{}{}{[Train1]}
\entry{}{}{}{[Train1]}
\end{nextentries}

% % Extension template
% \begin{nextentries}{Differential Geometry Notations, Part 3}{study:diffgeometry3}
% \entry{}{}{}{[Train1]}
% \entry{}{}{}{[Train1]}
% \entry{}{}{}{[Train1]}
% \entry{}{}{}{[Train1]}
% \entry{}{}{}{[Train1]}
% \entry{}{}{}{[Train1]}
% \entry{}{}{}{[Train1]}
% \entry{}{}{}{[Train1]}
% \entry{}{}{}{[Train1]}
% \entry{}{}{}{[Train1]}
% \entry{}{}{}{[Train1]}
% \entry{}{}{}{[Train1]}
% \entry{}{}{}{[Train1]}
% \entry{}{}{}{[Train1]}
% \entry{}{}{}{[Train1]}
% \end{nextentries}


\ednote{Complex named enttity: ``$U(⁢1)$ Chern-Simons gauge theory.''}
\newpage

\section{Other fields}
\subsection{Quantum Physics}
\ednote{Bra-ket notation}
\ednote{computer science, biology, chemistry...}
$\vdots$

\chapter{Discussion}

\chapter{Conclusion}

\end{document}